%------------------------------------------------------------------------------%
\documentclass[fleqn,utf8,aspectratio=169,12pt]{beamer}

\usepackage{gaal}
\usepackage{mathtools}

\title{Sistemas Lineares}

%------------------------------------------------------------------------------%
\begin{document}

\addtitle

%------------------------------------------------------------------------------%
\section{Sistemas Lineares}

%------------------------------------------------------------------------------%
\begin{frame}{Sistemas Lineares}
  \emph{Sistemas lineares} \\
  são conjuntos de equações lineares nas mesmas incógnitas \\
  que precisam ser resolvidos simultaneamente
\end{frame}

%------------------------------------------------------------------------------%
\begin{frame}{Interpretação geométrica}
Cada equação linear em duas incógnitas representa uma reta
\[
  \pause
  x + y = 10
\]
\[
  \pause
  2x + 3y = 24
\]

\pause
\medskip
Resolver o sistema significa encontrar um ponto que \\
pertence às duas retas simultaneamente
\end{frame}

%------------------------------------------------------------------------------%
\begin{frame}{Aplicação de Sistemas Lineares}
\onslide<+->{Em aplicações, podemos ter milhares ou milhões de incógnitas}

\onslide<+->{As equações podem representar:}
\begin{itemize}[<+->]
    \item distribuição de recursos
    \item redes de transporte
    \item circuitos elétricos
    \item modelos econômicos
    \item problemas de otimização
    \item estruturas mecânicas
    \item fenômenos físicos
    \item aproximações numéricas de equações diferenciais
\end{itemize}
\end{frame}

\end{document}

% =========================================================
% EXEMPLO CIRCUITOS
% =========================================================

\begin{frame}{Exemplo: circuitos elétricos}

Considere um circuito elétrico com várias correntes desconhecidas.

As leis de Kirchhoff fornecem relações entre as correntes.

Por exemplo:

\[
\begin{cases}
I_1+I_2-I_3=0,\\
2I_1+3I_2=10,\\
I_2+4I_3=8.
\end{cases}
\]

As incógnitas

\[
I_1,\ I_2,\ I_3
\]

podem ser determinadas resolvendo um sistema linear.

\medskip

Assim, o sistema não é o problema em si.

\begin{center}
\textbf{Ele é uma representação matemática do problema.}
\end{center}

\end{frame}

% =========================================================
% EXEMPLO PRODUÇÃO
% =========================================================

\begin{frame}{Exemplo: produção}

Uma empresa fabrica três produtos:

\[
A,\qquad B,\qquad C.
\]

Cada produto utiliza quantidades diferentes de três recursos.

Se conhecemos a quantidade total disponível de cada recurso,
podemos escrever equações relacionando as quantidades produzidas:

\[
\begin{cases}
a_{11}x+a_{12}y+a_{13}z=b_1,\\
a_{21}x+a_{22}y+a_{23}z=b_2,\\
a_{31}x+a_{32}y+a_{33}z=b_3.
\end{cases}
\]

Novamente aparece um sistema linear.

\end{frame}

% =========================================================
% A MATRIZ
% =========================================================

\begin{frame}{Uma nova forma de representar o sistema}

O sistema

\[
\begin{cases}
a_{11}x_1+\cdots+a_{1n}x_n=b_1,\\
a_{21}x_1+\cdots+a_{2n}x_n=b_2,\\
\vdots\\
a_{m1}x_1+\cdots+a_{mn}x_n=b_m
\end{cases}
\]

pode ser escrito como

\[
A\mathbf{x}=\mathbf{b},
\]

onde

\[
A=
\begin{pmatrix}
a_{11}&\cdots&a_{1n}\\
\vdots&&\vdots\\
a_{m1}&\cdots&a_{mn}
\end{pmatrix},
\]

\[
\mathbf{x}=
\begin{pmatrix}
x_1\\
\vdots\\
x_n
\end{pmatrix},
\qquad
\mathbf{b}=
\begin{pmatrix}
b_1\\
\vdots\\
b_m
\end{pmatrix}.
\]

\end{frame}

% =========================================================
% VANTAGEM
% =========================================================

\begin{frame}{Por que essa representação é importante?}

A notação

\[
A\mathbf{x}=\mathbf{b}
\]

permite separar:

\begin{itemize}
    \item $A$: informações sobre as relações entre as variáveis;
    \item $\mathbf{x}$: quantidades desconhecidas;
    \item $\mathbf{b}$: dados conhecidos do problema.
\end{itemize}

\medskip

Além disso, podemos desenvolver métodos sistemáticos para resolver
problemas com muitas incógnitas.

\medskip

Um desses métodos é o

\begin{center}
    \Large
    \textbf{escalonamento de matrizes}.
\end{center}

\end{frame}

% =========================================================
% MAS POR QUE LINEAR?
% =========================================================

\begin{frame}{Por que ``linear''?}

Uma equação linear possui incógnitas apenas na primeira potência
e não apresenta produtos entre incógnitas.

Por exemplo:

\[
2x+3y-z=5
\]

é linear.

Já

\[
xy+z=5
\]

não é linear.

Também não são lineares:

\[
x^2+y=4
\]

e

\[
\sin(x)+y=1.
\]

\medskip

A linearidade permite utilizar uma estrutura matemática
particularmente poderosa.

\end{frame}

% =========================================================
% GRANDE APLICAÇÃO
% =========================================================

\begin{frame}{Uma aplicação mais sofisticada}

Agora vamos considerar um problema diferente:

\begin{center}
    \Large
    Como determinar a temperatura em uma região?
\end{center}

Imagine uma placa metálica cuja temperatura é mantida fixa
em sua fronteira.

No interior da placa, a temperatura se distribui de acordo com
um princípio físico.

\end{frame}

% =========================================================
% LAPLACE
% =========================================================

\begin{frame}{A equação de Laplace}

Em regime estacionário, sem fontes ou sumidouros de calor,
a temperatura $u(x,y)$ pode ser modelada pela equação de Laplace:

\[
\frac{\partial^2u}{\partial x^2}
+
\frac{\partial^2u}{\partial y^2}
=0.
\]

Ou, de forma compacta,

\[
\Delta u=0.
\]

\medskip

Temos agora uma equação diferencial parcial.

\medskip

Como poderíamos utilizar um computador para obter uma
aproximação da solução?

\end{frame}

% =========================================================
% DISCRETIZAÇÃO
% =========================================================

\begin{frame}{Transformando o problema contínuo em um problema discreto}

Em vez de procurar a temperatura em todos os pontos da placa,
vamos considerar uma malha.

\medskip

Por exemplo:

\[
\begin{array}{ccccc}
\bullet & \bullet & \bullet & \bullet & \bullet\\
\bullet & \bullet & \bullet & \bullet & \bullet\\
\bullet & \bullet & \bullet & \bullet & \bullet\\
\bullet & \bullet & \bullet & \bullet & \bullet\\
\bullet & \bullet & \bullet & \bullet & \bullet
\end{array}
\]

Queremos aproximar a temperatura apenas nos pontos da malha.

\medskip

Chamaremos a temperatura no ponto $(x_i,y_j)$ de

\[
u_{ij}.
\]

\end{frame}

% =========================================================
% DIFERENÇAS FINITAS
% =========================================================

\begin{frame}{Aproximação por diferenças finitas}

Para uma malha uniforme, podemos aproximar as derivadas
segundas por diferenças finitas.

Temos aproximadamente:

\[
\frac{\partial^2u}{\partial x^2}
\approx
\frac{u_{i+1,j}-2u_{i,j}+u_{i-1,j}}{h^2}
\]

e

\[
\frac{\partial^2u}{\partial y^2}
\approx
\frac{u_{i,j+1}-2u_{i,j}+u_{i,j-1}}{h^2}.
\]

Substituindo na equação de Laplace:

\[
\frac{u_{i+1,j}-2u_{i,j}+u_{i-1,j}}{h^2}
+
\frac{u_{i,j+1}-2u_{i,j}+u_{i,j-1}}{h^2}
=0.
\]

\end{frame}

% =========================================================
% EQUAÇÃO DISCRETA
% =========================================================

\begin{frame}{A equação discreta}

Multiplicando por $h^2$:

\[
u_{i+1,j}-2u_{i,j}+u_{i-1,j}
+
u_{i,j+1}-2u_{i,j}+u_{i,j-1}
=0.
\]

Reorganizando:

\[
4u_{i,j}
=
u_{i+1,j}
+
u_{i-1,j}
+
u_{i,j+1}
+
u_{i,j-1}.
\]

Portanto,

\[
u_{i,j}
=
\frac{
u_{i+1,j}
+
u_{i-1,j}
+
u_{i,j+1}
+
u_{i,j-1}
}{4}.
\]

\medskip

A temperatura em cada ponto interior é aproximadamente a
\textbf{média das temperaturas dos quatro vizinhos}.

\end{frame}

% =========================================================
% PEQUENA MALHA
% =========================================================

\begin{frame}{Uma pequena malha}

Considere uma malha com quatro pontos interiores:

\[
\begin{array}{ccccc}
10 & 10 & 10 & 10 & 10\\
10 & u_1 & u_2 & 10 & 10\\
10 & u_3 & u_4 & 10 & 10\\
10 & 10 & 10 & 10 & 10
\end{array}
\]

Os valores $10$ representam temperaturas conhecidas na fronteira.

As incógnitas são

\[
u_1,\quad u_2,\quad u_3,\quad u_4.
\]

\end{frame}

% =========================================================
% EQUAÇÕES
% =========================================================

\begin{frame}{Obtendo as equações}

Para $u_1$:

\[
4u_1=10+u_2+10+u_3.
\]

Portanto,

\[
4u_1-u_2-u_3=20.
\]

Para $u_2$:

\[
4u_2=u_1+10+10+u_4,
\]

logo,

\[
-u_1+4u_2-u_4=20.
\]

\end{frame}

\begin{frame}{Obtendo as equações}

Para $u_3$:

\[
4u_3=10+u_4+10+u_1,
\]

portanto,

\[
-u_1+4u_3-u_4=20.
\]

Para $u_4$:

\[
4u_4=u_3+10+10+u_2,
\]

logo,

\[
-u_2-u_3+4u_4=20.
\]

\end{frame}

% =========================================================
% SISTEMA
% =========================================================

\begin{frame}{O sistema linear aparece naturalmente}

As quatro equações formam o sistema

\[
\begin{cases}
4u_1-u_2-u_3=20,\\
-u_1+4u_2-u_4=20,\\
-u_1+4u_3-u_4=20,\\
-u_2-u_3+4u_4=20.
\end{cases}
\]

Agora temos exatamente o tipo de problema que estamos estudando:

\[
A\mathbf{u}=\mathbf{b}.
\]

\end{frame}

% =========================================================
% MATRIZ
% =========================================================

\begin{frame}{Representação matricial}

O sistema pode ser escrito como

\[
\begin{pmatrix}
4&-1&-1&0\\
-1&4&0&-1\\
-1&0&4&-1\\
0&-1&-1&4
\end{pmatrix}
\begin{pmatrix}
u_1\\
u_2\\
u_3\\
u_4
\end{pmatrix}
=
\begin{pmatrix}
20\\
20\\
20\\
20
\end{pmatrix}.
\]

Ou simplesmente,

\[
A\mathbf{u}=\mathbf{b}.
\]

\medskip

O problema de determinar a distribuição de temperatura foi
transformado em um problema de \textbf{Álgebra Linear}.

\end{frame}

% =========================================================
% RESOLUÇÃO
% =========================================================

\begin{frame}{Resolvendo o sistema}

Por simetria, podemos esperar que

\[
u_1=u_2=u_3=u_4.
\]

Se

\[
u_1=u_2=u_3=u_4=u,
\]

a primeira equação fornece

\[
4u-u-u=20.
\]

Logo,

\[
2u=20
\]

e

\[
u=10.
\]

Portanto,

\[
u_1=u_2=u_3=u_4=10.
\]

\end{frame}

% =========================================================
% INTERPRETAÇÃO
% =========================================================

\begin{frame}{O que acabamos de fazer?}

Começamos com um problema físico:

\[
\text{distribuição de temperatura em uma placa}.
\]

\medskip

Utilizamos:

\[
\text{equação de Laplace}
\]

\[
\Downarrow
\]

\text{diferenças finitas}

\[
\Downarrow
\]

\text{equações lineares}

\[
\Downarrow
\]

\text{sistema linear}

\[
\Downarrow
\]

\text{matriz}

\[
\Downarrow
\]

\text{escalonamento / métodos numéricos}.

\end{frame}

% =========================================================
% MALHA MAIOR
% =========================================================

\begin{frame}{E se a malha for maior?}

Na nossa pequena malha tivemos apenas quatro incógnitas.

Mas imagine uma malha com:

\[
100\times100
\]

pontos interiores.

Teremos aproximadamente

\[
10\,000
\]

temperaturas desconhecidas.

O modelo produzirá um sistema com milhares de equações e
milhares de incógnitas.

\medskip

Ainda assim, sua estrutura será

\[
A\mathbf{u}=\mathbf{b}.
\]

\end{frame}

% =========================================================
% COMPUTAÇÃO
% =========================================================

\begin{frame}{E é aqui que a Álgebra Linear se torna computacional}

Para sistemas pequenos, podemos utilizar o escalonamento
manualmente.

Para sistemas com milhares ou milhões de incógnitas,
precisamos de algoritmos computacionais.

Alguns métodos importantes incluem:

\begin{itemize}
    \item eliminação de Gauss;
    \item fatoração $LU$;
    \item métodos iterativos;
    \item métodos específicos para matrizes esparsas.
\end{itemize}

\medskip

A Álgebra Linear fornece a linguagem e as ferramentas para
esses métodos.

\end{frame}

% =========================================================
% ESTRUTURA DA MATRIZ
% =========================================================

\begin{frame}{Uma característica importante}

Observe a matriz

\[
A=
\begin{pmatrix}
4&-1&-1&0\\
-1&4&0&-1\\
-1&0&4&-1\\
0&-1&-1&4
\end{pmatrix}.
\]

A maioria dos elementos é zero.

Uma matriz com muitos elementos iguais a zero é chamada de
\textbf{matriz esparsa}.

\medskip

Essa estrutura não é um detalhe:

\begin{center}
\Large
\textbf{ela pode ser explorada pelos algoritmos numéricos.}
\end{center}

\end{frame}

% =========================================================
% A IDEIA CENTRAL
% =========================================================

\begin{frame}{A ideia central}

Sistemas lineares não são apenas um tópico algébrico.

Eles aparecem quando transformamos problemas reais em modelos
matemáticos.

\medskip

Um mesmo objeto matemático,

\[
A\mathbf{x}=\mathbf{b},
\]

pode representar:

\begin{itemize}
    \item um circuito elétrico;
    \item um problema de produção;
    \item uma rede de transporte;
    \item uma estrutura mecânica;
    \item uma aproximação de uma equação diferencial;
    \item a distribuição de temperatura em uma placa.
\end{itemize}

\end{frame}

% =========================================================
% FECHAMENTO
% =========================================================

\begin{frame}{Por que estudar sistemas lineares?}

Porque aprender a resolver

\[
A\mathbf{x}=\mathbf{b}
\]

significa adquirir uma ferramenta que aparece em
diversas áreas da Matemática, da Ciência e da Engenharia.

\medskip

O escalonamento que fazemos com algumas linhas e colunas em uma
aula pode ser o mesmo princípio matemático utilizado para resolver
sistemas com milhares ou milhões de incógnitas.

\medskip

\begin{center}
\Large
\textbf{O sistema linear é uma ponte entre a modelagem e a solução.}
\end{center}

\end{frame}

%------------------------------------------------------------------------------%
% \framedgraphic{Gráfico de uma Função Real}{B_01_grafico_funcao_real}{}

%------------------------------------------------------------------------------%
\end{document}
%------------------------------------------------------------------------------%
